\section{Strategic Adaptation Diagnostics}
\label{app:strategic_adaptation_v18}

The award-layer screen is not invariant to strategic response.
This appendix records the diagnostic simulation that holds the
empirical environment fixed and asks how the participation
score behaves under simple evasive changes. The scenarios are
intentionally simple: they are not a full equilibrium model of
evasion, and they do not solve the strategic problem at its
root. Their purpose is to identify which kinds of adaptation
weaken a static binary rule more than others, and to inform the
operational recommendations in the main text.

\subsection{Simulation Scenarios}
\label{app:adaptation_scenarios_v18}

Five candidate adaptation strategies are tested against the
baseline frequent-loser rule. Rotation of cover bidders
distributes losing participation across a larger pool of
identities. Occasional wins introduce a small positive
win-share so that targeted firms do not satisfy the strict
$W_i=0$ filter. CNPJ splitting fragments a single firm's
participation across multiple legal identifiers. Threshold-aware
capping holds participation just below the median-plus-$1.5\times$IQR
cutoff so that the firm avoids the flag. The final scenario
combines capping with rotation.

Rotation leaves AUC essentially unchanged
($\valAUCAdaptRotation$ versus baseline $\valAUCBase$).
Occasional wins also leave AUC nearly unchanged
($\valAUCAdaptOccWins$), because the strict zero-win rule
removes those firms from the flagged set rather than
mis-ranking them. CNPJ splitting lowers AUC to
$\valAUCAdaptCNPJsplit$ by breaking the participation history
that the persistence signal depends on. Threshold-aware capping
is more damaging, lowering AUC to $\valAUCAdaptThreshCap$.
Combining capping with rotation lowers AUC to
$\valAUCAdaptCombined$.

\begin{table}[htbp]
\centering
\caption{Strategic-Adaptation Summary}
\label{tab:adaptation_summary_v18}
\footnotesize
\begin{adjustbox}{max width=\textwidth,center}
\begin{tabular}{p{4.4cm}cc}
\toprule
Scenario & AUC & Change vs baseline \\
\midrule
Baseline & $\valAUCBase$ & --- \\
Rotation & $\valAUCAdaptRotation$ & $\valAdaptRotationDeltaPP$ pp \\
Occasional wins & $\valAUCAdaptOccWins$ & $\valAdaptOccWinsDeltaPP$ pp \\
CNPJ splitting & $\valAUCAdaptCNPJsplit$ & $\valAdaptCNPJDeltaPP$ pp \\
Threshold capping & $\valAUCAdaptThreshCap$ & $\valAdaptCapDeltaPP$ pp \\
Capping plus rotation & $\valAUCAdaptCombined$ & $\valAdaptCombinedDeltaPP$ pp \\
\bottomrule
\end{tabular}
\end{adjustbox}
\end{table}

\subsection{Implementation Implications}
\label{app:adaptation_implementation_v18}

The hierarchy is informative for implementation. The screen is
not equally vulnerable to every adaptation strategy. Simple
rotation and occasional wins are absorbed by the persistence
signal; identity fragmentation begins to break the participation
history; behavior targeted at the static cutoff is the most
damaging. The implementation lessons are correspondingly
specific. The rule should be refreshed periodically as new
participation data arrive. Bunching just below the cutoff should
be monitored, because threshold-aware capping leaves a
detectable pile-up. Continuous ranks or capacity-constrained
top-$k$ queues should be preferred where they are operationally
feasible, since they are harder to game than a single published
cutoff. And the award-layer screen should be paired with
bid-layer forensics whenever bid files can be recovered, so
that strategic adaptation against the cheap layer does not also
succeed against the costly layer.

The simulation is a diagnostic, not a full equilibrium model of
evasion. Strategic response weakens a static binary rule more
than it invalidates the two-stage architecture.
